\chapter{Discussion}
\label{ch:discussion}

The \glsauto{ch4} yield data obtained throughout the experimental period is summarised in \Cref{tab:performance}, which reports the cumulative \glsauto{ch4} production, specific yield, and volumetric productivity for each of the four operational phases. The cumulative specific \glsauto{ch4} yield is calculated using \Cref{eq:cumYield}, while the daily specific \glsauto{ch4} production rate is obtained from \Cref{eq:dailyRate}. The volumetric productivity is derived from \Cref{eq:volRate}, and the phase averaged values are computed using \Cref{eq:avgMolRate,eq:avgVolRate,eq:avgDailyRate}. The overall yield for the entire experiment is given by \Cref{eq:overallYield}. These parameters are visualised in \Cref{graph:specific-accumulation-ch4-yield,graph:daily-yiled-ch4,graph:volumetric-productivity-ch4}, which show the evolution of \glsauto{ch4} yield and \glsauto{olr} across the four phases.

\begingroup
\tablesetup
\footnotesize
\setlength{\tabcolsep}{3pt}
\begin{ThreePartTable}
    \begin{xltabular}{\linewidth}{@{} c C c C C c c @{}}

        \caption{Summary of the performance of the experiment}
        \label{tab:performance} \\

        \toprule
        \thead{Phase} &
        \thead{Days} &
        \thead{\glsauto{olr} range\\(\unit{\gram\VS\per\litre\per\day})} &
        \thead{\glsauto{vs}\\(\unit{\gram})} &
        \thead{\glsauto{ch4}\\(\unit{\nml})} &
        \thead{Yield\\(\unit{\nml\per\gram\VSfed})} &
        \thead{Rate\\(\unit{\litre\per\reactorlitre\per\day})} \\
        \midrule
        \endfirsthead

        \multicolumn{7}{l}{\textit{\Cref{tab:performance} continued}} \\
        \toprule
        \thead{Phase} &
        \thead{Days} &
        \thead{\glsauto{olr} range\\(\unit{\gram\VS\per\litre\per\day})} &
        \thead{\glsauto{vs}\\(\unit{\gram})} &
        \thead{\glsauto{ch4}\\(\unit{\nml})} &
        \thead{Yield\\(\unit{\nml\per\gram\VSfed})} &
        \thead{Rate\\(\unit{\litre\per\reactorlitre\per\day})} \\
        \midrule
        \endhead

        \bottomrule
        \endlastfoot

        \csvreader[
            late after line=\\\addlinespace
        ]{data/processed/yield-avarage.csv}{}
        {\csvcoli &
         \numrange[parse-numbers=false]{\csvcolii}{\csvcoliii} &
         \numrange{\csvcoliv}{\csvcolv} &
         \num{\csvcolxi} &
         \num{\csvcolviii} &
         \num{\csvcolxii} &
         \num{\csvcolxiv}}

    \end{xltabular}
\end{ThreePartTable}
\endgroup

The overall cumulative \glsauto{ch4} production reached \qty{3276.1}{\nml}, corresponding to an overall yield of approximately \qty{28.96}{\nml\CH\per\gram\VSfed} (calculated using \Cref{eq:overallYield}). This value is considerably lower than the yields reported in the literature for similar substrates. \textcite{gebreegziabher2025} reviewed the \glsauto{ad} of biodiesel byproducts and reported \glsauto{ch4} yield enhancements ranging from \qtyrange{14}{226}{\percent} when \glsauto{cgl} is co-digested with various substrates. \textcite{tomczak2025} analysed data from multiple studies and concluded that \glsauto{ch4} content in biogas exceeds \qty{60}{\percent} when the \glsauto{cgl} fraction in the feedstock is kept below \qty{8}{\percent} (\glsauto{vpv}). In the present study, the \glsauto{cgl} fraction varied across phases, but the overall \glsauto{ch4} yield remained below the values typically reported for \glsauto{cgl} co-digestion. This discrepancy may be attributed to the single-stage configuration, the specific characteristics of the \glsauto{fw} and \glsauto{gl} used, and the process instability observed during the end of the co-digestion phase and the start of the \glsauto{gl}-only phase.

During the adaptation phase (days \numrange[parse-numbers=false]{0}{24}), the specific \glsauto{ch4} yield reached \qty{45.64}{\nml\CH\per\gram\VSfed}, the highest among all phases. However, this value must be interpreted with caution, as the cumulative \glsauto{vs} fed during this period was only \qty{16.628}{\gram}, and \glsauto{ch4} production was still in the early stages of development. The volumetric productivity was \qty{0.0304}{\litre\per\reactorlitre\per\day}. The daily yield profile in \Cref{graph:daily-yiled-ch4} shows a gradual increase in \glsauto{ch4} production during this phase, with a notable spike around day \num[parse-numbers=false]{8}, corresponding to the first addition of \glsauto{cgl} to the feed. This observation is consistent with \textcite{gebreegziabher2025}, who noted that \glsauto{cgl} is readily biodegradable and can stimulate methanogenic activity when introduced to an acclimated microbial community.

In the co-digestion phase (days \numrange[parse-numbers=false]{25}{52}), the \glsauto{olr} was increased from \qtyrange{1.4}{3.16}{\gram\VS\per\litre\per\day}, and the cumulative \glsauto{ch4} production reached \qty{1249.9}{\nml}. The specific yield decreased to \qty{19.71}{\nml\CH\per\gram\VSfed}, while the volumetric productivity increased to \qty{0.0446}{\litre\per\reactorlitre\per\day}, the highest of all phases. This inverse relationship between specific yield and volumetric productivity reflects the higher organic loading rate applied, which promoted greater \glsauto{ch4} production per reactor volume but reduced the efficiency of substrate conversion to \glsauto{ch4}. The daily yield profile in \Cref{graph:daily-yiled-ch4} exhibits fluctuations during this phase, this pattern is typical of readily degradable substrates and has been reported by \textcite{tomczak2025} and \textcite{srimachai2025}. \textcite{srimachai2025} investigated the co-digestion of canned seafood wastewater with \glsauto{gl} waste and reported \glsauto{ch4} yields of \qty{345}{\mL\CH\per\gram\COD} at a mixing ratio of \qty{96.4}{\percent} (\glsauto{vpv}). Although the substrate characteristics differ, the general trend of enhanced \glsauto{ch4} production during co-digestion is consistent with the present findings.

The \glsauto{cgl} only phase (days \numrange[parse-numbers=false]{53}{81}) exhibited a specific yield of \qty{37.37}{\nml\CH\per\gram\VSfed} and a volumetric productivity of \qty{0.0303}{\litre\per\reactorlitre\per\day}, with a cumulative \glsauto{ch4} production of \qty{877.9}{\nml}. The \glsauto{olr} ranged from \qtyrange{0.42}{1.37}{\gram\VS\per\litre\per\day}. Despite the absence of \glsauto{fw} in the feed, \glsauto{ch4} production continued, indicating that the microbial community had adapted to \glsauto{gl} as a sole carbon source. However, the specific yield was lower than during the adaptation phase, suggesting that \glsauto{gl} alone may not provide an optimal nutrient balance for methanogenesis. \textcite{tomczak2025} emphasised that \glsauto{cgl} is an ideal co-substrate but can cause inhibition when used alone or at high concentrations due to the accumulation of \glsplauto{vfa} and the presence of impurities such as methanol, soaps, and heavy metals. The daily yield profile in \Cref{graph:daily-yiled-ch4} shows a declining trend during this phase, with a notable outlier peak at day \num[parse-numbers=false]{53}, corresponding to the transition from co-digestion to \glsauto{gl}-only feeding. This anomaly, also visible in \Cref{graph:volumetric-productivity-ch4}, was caused by problems during the change.

The post-feeding phase (days \numrange[parse-numbers=false]{82}{97}) produced \qty{389.4}{\nml} of \glsauto{ch4}, with a specific yield of \qty{40.60}{\nml\CH\per\gram\VSfed} and a volumetric productivity of \qty{0.0245}{\litre\per\reactorlitre\per\day}. The \glsauto{olr} got to zero during this period, and the \glsauto{ch4} produced can be attributed to the slow degradation of residual organic matter and microbial decay products. The daily yield profile in \Cref{graph:daily-yiled-ch4} shows a gradual decline, confirming that the reactor was no longer being actively fed and that the remaining substrate was limited.

\begingroup

\definecolor{phaseAdaptation}{HTML}{D6E4F0}
\definecolor{phaseCodigestion}{HTML}{D5EAD5}
\definecolor{phaseCGL}{HTML}{FBE5C9}
\definecolor{phasePost}{HTML}{F8D5D5}

\pgfplotstableread[col sep=comma]{data/processed/yield.csv}{\phasedata}%

\begin{figure}[H]
    \centering
    \begin{tikzpicture}
        % ---------- LEFT axis: Ph ----------
        \begin{axis}[
            width=\textwidth, height=0.6\textwidth,
            xlabel={Time (\unit{\day})},
            ylabel={\unit{\CH\milli\litre\per\gram\VSfed}},
            axis y line*=left,
            enlarge x limits=false,
            xmin=0, xmax=97,
            enlarge x limits=false,
            set layers=axis on top,
            legend style={
                at={(0.98, 0.8)},
                anchor=east,
                draw=black,
                fill=white,
                fill opacity=0.9,
            },
        ]
            % 1. Data first — this is what lets pgfplots compute the limits
            \addplot[very thick, mark=*, mark size=1.2pt, color=blue!60!black]
                table[x index=0, y index=6] {\phasedata};
            \addlegendentry{\glsauto{ch4}}

            \addlegendimage{color=red!60!black, very thick, mark=*, mark size=1.2pt}
            \addlegendentry{OLR}

            % 2. Shading + labels deferred until limits are known
            \pgfplotsextra{%
                \pgfonlayer{axis background}
                    \fill[phaseAdaptation]
                    (axis cs:\pgfkeysvalueof{/pgfplots/xmin},\pgfkeysvalueof{/pgfplots/ymin})
                    rectangle (axis cs:24.5,\pgfkeysvalueof{/pgfplots/ymax});
                    \fill[phaseCodigestion]
                    (axis cs:24.5,\pgfkeysvalueof{/pgfplots/ymin})
                    rectangle (axis cs:52.5,\pgfkeysvalueof{/pgfplots/ymax});
                    \fill[phaseCGL]
                    (axis cs:52.5,\pgfkeysvalueof{/pgfplots/ymin})
                    rectangle (axis cs:81.5,\pgfkeysvalueof{/pgfplots/ymax});
                    \fill[phasePost]
                    (axis cs:81.5,\pgfkeysvalueof{/pgfplots/ymin})
                    rectangle (axis cs:\pgfkeysvalueof{/pgfplots/xmax},\pgfkeysvalueof{/pgfplots/ymax});
                \endpgfonlayer
                \node[font=\scriptsize, text=black!70, anchor=north]
                    at (axis cs:12, \pgfkeysvalueof{/pgfplots/ymax}) {Adaptation};
                \node[font=\scriptsize, text=black!70, anchor=north]
                    at (axis cs:38, \pgfkeysvalueof{/pgfplots/ymax}) {Co-digestion};
                \node[font=\scriptsize, text=black!70, anchor=north]
                    at (axis cs:67, \pgfkeysvalueof{/pgfplots/ymax}) {C-GL only};
                \node[font=\scriptsize, text=black!70, anchor=north]
                    at (axis cs:89, \pgfkeysvalueof{/pgfplots/ymax}) {Post-feeding};
            }
        \end{axis}

        % ---------- RIGHT axis: Ratio ----------
        \begin{axis}[
            width=\textwidth, height=0.6\textwidth,
            axis y line*=right,       % only right spine + ticks
            axis x line=none,         % hide duplicate x-axis
            axis background/.style={fill=none},
            xmin=0, xmax=97,
            enlarge x limits=false,
            set layers=axis on top,
            ylabel={ORL (\unit{\kilo\gram\VS\per\cubic\metre\per\day})},
            % no legend here — we put it on the left axis
            ]
            \addplot[very thick, mark=*, mark size=1.2pt, color=red!60!black]
                table[x index=0, y index=3] {\phasedata};
        \end{axis}
    \end{tikzpicture}
    \caption{Cumulative specific \glsentrysymbolauto{ch4} yield}
    \label{graph:specific-accumulation-ch4-yield}
\end{figure}

\begin{figure}[H]
    \centering
    \begin{tikzpicture}
        \begin{axis}[
            width=\textwidth, height=0.6\textwidth,
            xlabel={Time (\unit{\day})},
            ylabel={\unit{\CH\milli\litre\per\gram\VSfed\per\day}},
            axis y line*=left,
            enlarge x limits=false,
            xmin=0, xmax=97,
            enlarge x limits=false,
            set layers=axis on top,
            legend style={
                at={(0.98, 0.8)},
                anchor=east,
                draw=black,
                fill=white,
                fill opacity=0.9,
            },
        ]
            % 1. Data first — this is what lets pgfplots compute the limits
            \addplot[very thick, mark=*, mark size=1.2pt, color=blue!60!black]
                table[x index=0, y index=9] {\phasedata};
            \addlegendentry{\glsauto{ch4}}

            \addlegendimage{color=red!60!black, very thick, mark=*, mark size=1.2pt}
            \addlegendentry{OLR}


            % 2. Shading + labels deferred until limits are known
            \pgfplotsextra{%
                \pgfonlayer{axis background}
                    \fill[phaseAdaptation]
                    (axis cs:\pgfkeysvalueof{/pgfplots/xmin},\pgfkeysvalueof{/pgfplots/ymin})
                    rectangle (axis cs:24.5,\pgfkeysvalueof{/pgfplots/ymax});
                    \fill[phaseCodigestion]
                    (axis cs:24.5,\pgfkeysvalueof{/pgfplots/ymin})
                    rectangle (axis cs:52.5,\pgfkeysvalueof{/pgfplots/ymax});
                    \fill[phaseCGL]
                    (axis cs:52.5,\pgfkeysvalueof{/pgfplots/ymin})
                    rectangle (axis cs:81.5,\pgfkeysvalueof{/pgfplots/ymax});
                    \fill[phasePost]
                    (axis cs:81.5,\pgfkeysvalueof{/pgfplots/ymin})
                    rectangle (axis cs:\pgfkeysvalueof{/pgfplots/xmax},\pgfkeysvalueof{/pgfplots/ymax});
                \endpgfonlayer
                \node[font=\scriptsize, text=black!70, anchor=north]
                    at (axis cs:12, \pgfkeysvalueof{/pgfplots/ymax}) {Adaptation};
                \node[font=\scriptsize, text=black!70, anchor=north]
                    at (axis cs:38, \pgfkeysvalueof{/pgfplots/ymax}) {Co-digestion};
                \node[font=\scriptsize, text=black!70, anchor=north]
                    at (axis cs:67, \pgfkeysvalueof{/pgfplots/ymax}) {C-GL only};
                \node[font=\scriptsize, text=black!70, anchor=north]
                    at (axis cs:89, \pgfkeysvalueof{/pgfplots/ymax}) {Post-feeding};
            }
        \end{axis}

        % ---------- RIGHT axis: Ratio ----------
        \begin{axis}[
            width=\textwidth, height=0.6\textwidth,
            axis y line*=right,       % only right spine + ticks
            axis x line=none,         % hide duplicate x-axis
            axis background/.style={fill=none},
            xmin=0, xmax=97,
            enlarge x limits=false,
            set layers=axis on top,
            ylabel={ORL (\unit{\kilo\gram\VS\per\cubic\metre\per\day})},
            % no legend here — we put it on the left axis
            ]
            \addplot[very thick, mark=*, mark size=1.2pt, color=red!60!black]
                table[x index=0, y index=3] {\phasedata};
        \end{axis}
    \end{tikzpicture}
    \caption{Daily yield of \glsentrysymbolauto{ch4}}
    \label{graph:daily-yiled-ch4}
\end{figure}

\begin{figure}[H]
    \centering
    \begin{tikzpicture}
        \begin{axis}[
            width=\textwidth, height=0.6\textwidth,
            xlabel={Time (\unit{\day})},
            ylabel={\unit{\CH\litre\per\litre\per\day}},
            axis y line*=left,
            enlarge x limits=false,
            xmin=0, xmax=97,
            enlarge x limits=false,
            set layers=axis on top,
            legend style={
                at={(0.98, 0.8)},
                anchor=east,
                draw=black,
                fill=white,
                fill opacity=0.9,
            },
        ]
            % 1. Data first — this is what lets pgfplots compute the limits
            \addplot[very thick, mark=*, mark size=1.2pt, color=blue!60!black]
                table[x index=0, y index=8] {\phasedata};
            \addlegendentry{\glsauto{ch4}}

            \addlegendimage{color=red!60!black, very thick, mark=*, mark size=1.2pt}
            \addlegendentry{OLR}


            % 2. Shading + labels deferred until limits are known
            \pgfplotsextra{%
                \pgfonlayer{axis background}
                    \fill[phaseAdaptation]
                    (axis cs:\pgfkeysvalueof{/pgfplots/xmin},\pgfkeysvalueof{/pgfplots/ymin})
                    rectangle (axis cs:24.5,\pgfkeysvalueof{/pgfplots/ymax});
                    \fill[phaseCodigestion]
                    (axis cs:24.5,\pgfkeysvalueof{/pgfplots/ymin})
                    rectangle (axis cs:52.5,\pgfkeysvalueof{/pgfplots/ymax});
                    \fill[phaseCGL]
                    (axis cs:52.5,\pgfkeysvalueof{/pgfplots/ymin})
                    rectangle (axis cs:81.5,\pgfkeysvalueof{/pgfplots/ymax});
                    \fill[phasePost]
                    (axis cs:81.5,\pgfkeysvalueof{/pgfplots/ymin})
                    rectangle (axis cs:\pgfkeysvalueof{/pgfplots/xmax},\pgfkeysvalueof{/pgfplots/ymax});
                \endpgfonlayer
                \node[font=\scriptsize, text=black!70, anchor=north]
                    at (axis cs:12, \pgfkeysvalueof{/pgfplots/ymax}) {Adaptation};
                \node[font=\scriptsize, text=black!70, anchor=north]
                    at (axis cs:38, \pgfkeysvalueof{/pgfplots/ymax}) {Co-digestion};
                \node[font=\scriptsize, text=black!70, anchor=north]
                    at (axis cs:67, \pgfkeysvalueof{/pgfplots/ymax}) {C-GL only};
                \node[font=\scriptsize, text=black!70, anchor=north]
                    at (axis cs:89, \pgfkeysvalueof{/pgfplots/ymax}) {Post-feeding};
            }
        \end{axis}

        % ---------- RIGHT axis: Ratio ----------
        \begin{axis}[
            width=\textwidth, height=0.6\textwidth,
            axis y line*=right,       % only right spine + ticks
            axis x line=none,         % hide duplicate x-axis
            axis background/.style={fill=none},
            xmin=0, xmax=97,
            enlarge x limits=false,
            set layers=axis on top,
            ylabel={ORL (\unit{\kilo\gram\VS\per\cubic\metre\per\day})},
            % no legend here — we put it on the left axis
            ]
            \addplot[very thick, mark=*, mark size=1.2pt, color=red!60!black]
                table[x index=0, y index=3] {\phasedata};
        \end{axis}
    \end{tikzpicture}
    \caption{Volumetric productivity of \glsentrysymbolauto{ch4}}
    \label{graph:volumetric-productivity-ch4}
\end{figure}

\endgroup

\begin{equation}
    Q = \frac{V_r}{\mathrm{HRT}} \quad \left[\unit{\cubic\metre\per\day}\right]
    \label{eq:Q}
\end{equation}

\begin{equation}
    m_{\mathrm{VS,day}} = \mathrm{OLR} \cdot V_r \quad \left[\unit{\gram\per\day}\right]
    \label{eq:mVS_day}
\end{equation}

\begin{equation}
    M_{\mathrm{VS}}(n) = \sum_{i=0}^{n} m_{\mathrm{VS,day}}(i) \quad \left[\unit{\gram}\right]
    \label{eq:cumVS}
\end{equation}

\begin{equation}
    Y(n) = \frac{V_{\mathrm{CH_4}}(n)}{M_{\mathrm{VS}}(n)} \quad \left[\unit{\nml\CH\per\gram\VSfed}\right]
    \label{eq:cumYield}
\end{equation}

\begin{equation}
    r_{\mathrm{CH_4}}(n) = \frac{F_{\mathrm{CH_4}}(n)}{m_{\mathrm{VS,day}}(n)} \quad \left[\unit{\milli\litre\CH\per\gram\VSfed\per\day}\right]
    \label{eq:dailyRate}
\end{equation}

\begin{equation}
    \dot{n}_{\mathrm{CH_4}}(n) = \frac{F_{\mathrm{CH_4}}(n)}{V_m \cdot V_r} \quad \left[\unit{\mol\CH\per\cubic\metre\per\day}\right]
    \label{eq:molRate}
\end{equation}

\begin{equation}
    \dot{V}_{\mathrm{CH_4}}(n) = \frac{F_{\mathrm{CH_4}}(n)}{1000 \cdot V_r} \quad \left[\unit{\litre\CH\per\reactorlitre\per\day}\right]
    \label{eq:volRate}
\end{equation}

\begin{equation}
    \Delta V_{\mathrm{CH_4}}(a,b) = V_{\mathrm{CH_4}}(b) - V_{\mathrm{CH_4}}(a-1) \quad \left[\unit{\nml}\right]
    \label{eq:deltaVCH4}
\end{equation}

\begin{equation}
    \Delta M_{\mathrm{VS}}(a,b) = M_{\mathrm{VS}}(b) - M_{\mathrm{VS}}(a-1) \quad \left[\unit{\gram}\right]
    \label{eq:deltaMVS}
\end{equation}

\begin{equation}
    Y(a,b) = \frac{\Delta V_{\mathrm{CH_4}}(a,b)}{\Delta M_{\mathrm{VS}}(a,b)} \quad \left[\unit{\nml\CH\per\gram\VSfed}\right]
    \label{eq:phaseYield}
\end{equation}

\begin{equation}
    \overline{\dot{n}_{\mathrm{CH_4}}}(a,b) = \frac{1}{b - a + 1} \sum_{n=a}^{b} \dot{n}_{\mathrm{CH_4}}(n) \quad \left[\unit{\mol\CH\per\cubic\metre\per\day}\right]
    \label{eq:avgMolRate}
\end{equation}

\begin{equation}
    \overline{\dot{V}_{\mathrm{CH_4}}}(a,b) = \frac{1}{b - a + 1} \sum_{n=a}^{b} \dot{V}_{\mathrm{CH_4}}(n) \quad \left[\unit{\litre\CH\per\reactorlitre\per\day}\right]
    \label{eq:avgVolRate}
\end{equation}

\begin{equation}
    \overline{r_{\mathrm{CH_4}}}(a,b) = \frac{1}{b - a + 1} \sum_{n=a}^{b} r_{\mathrm{CH_4}}(n) \quad \left[\unit{\milli\litre\CH\per\gram\VSfed\per\day}\right]
    \label{eq:avgDailyRate}
\end{equation}

\begin{equation}
    Y_{\text{overall}} = \frac{V_{\mathrm{CH_4}}(N_f)}{M_{\mathrm{VS}}(N_f)} \quad \left[\unit{\nml\CH\per\gram\VSfed}\right]
    \label{eq:overallYield}
\end{equation}

Comparisons with the literature reveal that the \glsauto{ch4} yields obtained in this experiment are lower than those reported for similar co-digestion systems. \textcite{aguiar2026} investigated the co-digestion of sugarcane vinasse and cattle manure in a semi-continuous stirred-tank reactor and reported volumetric \glsauto{ch4} productivities of \qtyrange{3981}{5269}{\nml\CH\per\litre\per\day} for mixing ratios of 5:1 and 1:1 (v/m), respectively, at \glsplauto{olr} of \qtyrange{1.7}{4.1}{\kg\VS\per\cubic\metre\per\day}. These values are two orders of magnitude higher than the maximum volumetric productivity of \qty{44.6}{\nml\CH\per\litre\per\day} obtained in the present study. The large difference is primarily due to the much higher \glsauto{olr} applied by \textcite{aguiar2026} and the different substrate characteristics. \textcite{aguiar2026} also reported that the addition of \glsauto{gl} as a supplementary organic substrate compromised system stability, which is consistent with the \glsauto{vfa} accumulation and pH fluctuations observed in the present study during the co-digestion phase.

\textcite{longanesi2016} reported \glsauto{ch4} yields from grape pomace digestion ranging from \qty{20.4}{\NL\CH\per\kg\VS} in batch mode to \qty{160}{\NL\CH\per\kg\VS} in continuous mode at an \glsauto{olr} of \qty{0.4}{\g\VS\per\litre\per\day}. The continuous-mode yield of \qty{160}{\NL\CH\per\kg\VS} is equivalent to \qty{160}{\nml\CH\per\gram\VS}, which is higher than the overall yield of \qty{28.96}{\nml\CH\per\gram\VSfed} obtained in this experiment. However, it should be noted that \textcite{longanesi2016} used grape pomace as the sole substrate, whereas the present experiment used a mixture of \glsauto{fw} and \glsauto{cgl}. \textcite{longanesi2016} also reported that co-digestion of grape pomace with maize silage improved the \glsauto{ch4} yield by \qty{113.5}{\NL\CH\per\kg\VS}, highlighting the benefits of co-digestion for balancing the \glsauto{cn_ratio} and diluting inhibitory compounds.

The lower \glsauto{ch4} yields obtained in this study can be attributed to several factors. First, the single-stage configuration may have limited the hydrolysis and acidogenesis of the \glsauto{fw}, as the optimal conditions for these processes differ from those for methanogenesis. \textcite{srimachai2025} demonstrated that a two-stage system (STR - tubular reactor) can enhance \glsauto{ch4} production by separating the hydrogen producing and \glsauto{ch4}-producing phases. Second, the \glsauto{cgl} impurities may have inhibited methanogenic activity. \textcite{anand2012} reported that \glsauto{cgl} from biodiesel production contains impurities such as methanol, soaps, and free fatty acids that can inhibit microbial growth and fermentation. Third, the \glsauto{vfa_alk} fluctuated from \numrange{0.3}{1.2}, indicating potential process instability. \textcite{ribas2007} emphasised the importance of accurate \glsauto{vfa} and \glsauto{alk} measurements for timely detection of process imbalance. Fourth, the fungal contamination observed during the cleaning of the reactor may have competed with methanogens for substrate, potentially reducing \glsauto{ch4} production. \textcite{degueurce2020} reported that storage of \glsauto{fw} can lead to the growth of fungi and other microorganisms that consume organic matter and produce inhibitory compounds.

Despite the lower yields, the results demonstrate that \glsauto{ch4} production was maintained throughout the experimental period, confirming the technical feasibility of co-digesting \glsauto{fw} and \glsauto{cgl} in a single-stage system. The specific \glsauto{ch4} yield during the co-digestion phase (\qty{19.71}{\nml\CH\per\gram\VSfed}) was lower than during the adaptation phase, but the volumetric productivity was higher, indicating that the reactor was able to handle increased \glsplauto{olr}. These findings are consistent with \textcite{gebreegziabher2025}, who reported that co-digestion of \glsauto{cgl} with nitrogen-rich substrates such as \glsauto{fw} can enhance biogas production by balancing the \glsauto{cn_ratio} and providing essential nutrients. However, the accumulation of \glsplauto{vfa} and the decrease in pH observed during the co-digestion phase highlight the need for careful monitoring and control of the \glsauto{olr}. \textcite{tomczak2025} recommended that the \glsauto{cgl} content in the feedstock should not exceed \qty{8}{\percent} (\glsauto{vpv}) to avoid inhibition, and that the \glsauto{olr} should be adjusted based on the buffering capacity of the system.
